% ============================================================
% Chapter 7: Cryptographic Mechanisms and Public Key Infrastructure
% Language: Greek — edit freely; regenerate from src/content if preferred.
% ============================================================
\chapter{Κρυπτογραφικοί Μηχανισμοί και Υποδομή Δημόσιου Κλειδιού}\label{ch:7}
\chaptermeta{Συμμετρική και ασύμμετρη κρυπτογράφηση · Συναρτήσεις κατακερματισμού · Ψηφιακές υπογραφές · ΥΔΚ και X.509 · TLS 1.3 · Διαχείριση κλειδιών}{Επίπεδο: Μεσαίο επίπεδο · Εφαρμοσμένη κρυπτογραφία}{Χρόνος μελέτης: 9–11 ώρες}

Η κρυπτογραφία αποτελεί το μαθηματικό θεμέλιο της εμπιστευτικότητας, της ακεραιότητας και της αυθεντικότητας στα ψηφιακά συστήματα. Το Κεφάλαιο 6 ολοκληρώθηκε με επιθέσεις που υποκλέπτουν και αλλοιώνουν την κίνηση· το παρόν κεφάλαιο παρέχει τα εργαλεία που καθιστούν άχρηστη μια τέτοια υποκλοπή. Στόχος δεν είναι η παραγωγή των υποκείμενων μαθηματικών, αλλά η ακριβής κατανόηση του τι εγγυάται κάθε κρυπτογραφικό δομικό στοιχείο, πώς συνδυάζονται τα στοιχεία αυτά και πού αποτυγχάνουν τα πραγματικά συστήματα.

Το κεφάλαιο αντιπαραβάλλει αρχικά τη συμμετρική με την ασύμμετρη κρυπτογράφηση και εξηγεί τις συναρτήσεις κατακερματισμού και τις ψηφιακές υπογραφές. Στη συνέχεια δείχνει πώς η υποδομή δημόσιου κλειδιού (ΥΔΚ ή PKI) συνδέει κλειδιά με ταυτότητες μέσω πιστοποιητικών και αναλύει τη χειραψία του TLS 1.3, που προστατεύει σήμερα το μεγαλύτερο μέρος της κίνησης στο διαδίκτυο. Κλείνει με τη διαχείριση κλειδιών —το πεδίο όπου, στην πράξη, σημειώνονται οι περισσότερες κρυπτογραφικές αποτυχίες.

\begin{outcomesbox}{Μαθησιακά αποτελέσματα}
Μετά τη μελέτη του κεφαλαίου θα πρέπει να μπορείτε:
\begin{itemize}
  \item Να εξηγείτε τη διαφορά μεταξύ συμμετρικής και ασύμμετρης κρυπτογράφησης και γιατί τα πρακτικά συστήματα τις συνδυάζουν.
  \item Να διατυπώνετε τις ιδιότητες ασφάλειας των κρυπτογραφικών συναρτήσεων κατακερματισμού και να διακρίνετε τον κατακερματισμό από την κρυπτογράφηση.
  \item Να περιγράφετε πώς οι ψηφιακές υπογραφές παρέχουν ακεραιότητα, αυθεντικότητα και (υπό προϋποθέσεις) μη αποποίηση.
  \item Να εξηγείτε τον ρόλο των αρχών πιστοποίησης, των αρχών καταχώρισης, των αλυσίδων πιστοποιητικών και της ανάκλησης σε μια ΥΔΚ.
  \item Να σκιαγραφείτε τη χειραψία του TLS 1.3 και να εντοπίζετε συνήθεις αστοχίες στη διαχείριση κλειδιών.
\end{itemize}
\end{outcomesbox}

\section{Συμμετρική και ασύμμετρη κρυπτογράφηση}\label{sec:7.1}
Στη \textbf{συμμετρική κρυπτογράφηση} χρησιμοποιείται το ίδιο μυστικό κλειδί για την κρυπτογράφηση και την αποκρυπτογράφηση. Σύγχρονοι συμμετρικοί αλγόριθμοι, όπως οι \textbf{AES} και \textbf{ChaCha20}, είναι εξαιρετικά γρήγοροι και, με κλειδιά 128 ή 256 bit, θεωρούνται ασφαλείς απέναντι σε όλες τις γνωστές πρακτικές επιθέσεις. Η αδυναμία τους είναι οργανωτική: τα δύο μέρη πρέπει να μοιράζονται ήδη το κλειδί, ενώ ένα δίκτυο \emph{n} συμμετεχόντων θα χρειαζόταν περίπου n²/2 διαφορετικά κλειδιά. Οι συμμετρικοί αλγόριθμοι πρέπει επίσης να χρησιμοποιούνται με κατάλληλο \emph{τρόπο λειτουργίας}· οι τρόποι με ενσωματωμένη αυθεντικοποίηση, όπως ο AES-GCM, προστατεύουν τόσο την εμπιστευτικότητα όσο και την ακεραιότητα.

Η \textbf{ασύμμετρη κρυπτογράφηση} (κρυπτογράφηση δημόσιου κλειδιού) χρησιμοποιεί ένα ζεύγος μαθηματικά συσχετισμένων κλειδιών. Το \textbf{δημόσιο κλειδί} μπορεί να διανέμεται ελεύθερα, ενώ το \textbf{ιδιωτικό κλειδί} παραμένει μυστικό. Δεδομένα που κρυπτογραφούνται με το δημόσιο κλειδί αποκρυπτογραφούνται μόνο με το ιδιωτικό. Αλγόριθμοι όπως ο \textbf{RSA} και η \textbf{κρυπτογραφία ελλειπτικών καμπυλών (ECC)} επιλύουν το πρόβλημα της διανομής κλειδιών, είναι όμως χιλιάδες φορές πιο αργοί από τους συμμετρικούς. Γι' αυτό τα πρακτικά συστήματα χρησιμοποιούν \textbf{υβριδική κρυπτογράφηση}: οι ασύμμετρες τεχνικές δημιουργούν ή προστατεύουν ένα τυχαίο συμμετρικό \emph{κλειδί συνεδρίας} και αυτό το κλειδί κρυπτογραφεί στη συνέχεια τον κύριο όγκο των δεδομένων. Κάθε σύνδεση TLS —και, όπως έδειξε το Κεφάλαιο 4, κάθε σύγχρονη επίθεση λυτρισμικού— ακολουθεί αυτό το πρότυπο.

\section{Κρυπτογραφικές συναρτήσεις κατακερματισμού}\label{sec:7.2}
Μια \textbf{κρυπτογραφική συνάρτηση κατακερματισμού} αντιστοιχίζει είσοδο οποιουδήποτε μήκους σε έξοδο σταθερού μήκους, τη \emph{σύνοψη} (digest). Η SHA-256, για παράδειγμα, παράγει πάντοτε 256 bit. Τρεις ιδιότητες καθιστούν μια συνάρτηση κατακερματισμού κατάλληλη για εφαρμογές ασφάλειας. \textbf{Αντίσταση προεικόνας}: δεδομένης μιας σύνοψης, είναι υπολογιστικά ανέφικτο να βρεθεί είσοδος που την παράγει. \textbf{Αντίσταση δεύτερης προεικόνας}: δεδομένης μιας εισόδου, είναι ανέφικτο να βρεθεί διαφορετική είσοδος με την ίδια σύνοψη. \textbf{Αντίσταση συγκρούσεων}: είναι ανέφικτο να βρεθούν \emph{οποιεσδήποτε} δύο είσοδοι με την ίδια σύνοψη. Οι MD5 και SHA-1 έχουν απολέσει την αντίσταση συγκρούσεων και δεν πρέπει πλέον να χρησιμοποιούνται σε υπογραφές· οι SHA-256, SHA-3 και BLAKE2 παραμένουν ασφαλείς.

Ο κατακερματισμός δεν είναι κρυπτογράφηση: δεν υπάρχει κλειδί ούτε τρόπος «αποκρυπτογράφησης» μιας σύνοψης. Οι συνόψεις χρησιμοποιούνται για την επαλήθευση της ακεραιότητας αρχείων, για την ταυτοποίηση αποδεικτικών στοιχείων στην εγκληματολογική ανάλυση (Κεφάλαιο 12) και ως δομικά στοιχεία των υπογραφών. Σε συνδυασμό με ένα μυστικό κλειδί σχηματίζουν έναν κώδικα \textbf{HMAC}, ο οποίος αποδεικνύει ότι ένα μήνυμα προέρχεται από κάποιον που γνωρίζει το κλειδί. Για την αποθήκευση κωδικών πρόσβασης οι συνήθεις γρήγορες συναρτήσεις είναι ακατάλληλες, επειδή οι επιτιθέμενοι μπορούν να δοκιμάζουν δισεκατομμύρια μαντεψιές ανά δευτερόλεπτο. Ειδικές \textbf{συναρτήσεις κατακερματισμού κωδικών}, όπως οι Argon2, scrypt ή bcrypt, είναι σκόπιμα αργές και χρησιμοποιούν ένα μοναδικό τυχαίο \emph{αλάτι} (salt) για κάθε κωδικό.

\section{Ψηφιακές υπογραφές και μη αποποίηση}\label{sec:7.3}
Μια \textbf{ψηφιακή υπογραφή} αντιστρέφει τους ρόλους του ζεύγους κλειδιών. Ο υπογράφων υπολογίζει τη σύνοψη του μηνύματος και τη μετασχηματίζει με το \emph{ιδιωτικό} του κλειδί· οποιοσδήποτε μπορεί να επαληθεύσει το αποτέλεσμα με το αντίστοιχο \emph{δημόσιο} κλειδί. Μια έγκυρη υπογραφή αποδεικνύει τρία πράγματα: ότι το μήνυμα δεν έχει αλλοιωθεί από τη στιγμή της υπογραφής (\textbf{ακεραιότητα}), ότι υπογράφηκε από τον κάτοχο του ιδιωτικού κλειδιού (\textbf{αυθεντικότητα}) και ότι ο υπογράφων δεν μπορεί εύκολα να ισχυριστεί αργότερα το αντίθετο (\textbf{μη αποποίηση}). Συνήθεις αλγόριθμοι είναι οι RSA-PSS, ECDSA και Ed25519.

Η μη αποποίηση, ωστόσο, είναι τόσο ισχυρή όσο η διαδικασία που την περιβάλλει. Αν το ιδιωτικό κλειδί ήταν αποθηκευμένο χωρίς προστασία σε κοινόχρηστο υπολογιστή, ο υπογράφων μπορεί εύλογα να υποστηρίξει ότι το χρησιμοποίησε κάποιος άλλος. Γι' αυτό νομικά πλαίσια όπως ο Κανονισμός \textbf{eIDAS} της ΕΕ διακρίνουν τις απλές ηλεκτρονικές υπογραφές από τις \emph{εγκεκριμένες} ηλεκτρονικές υπογραφές, οι οποίες απαιτούν εγκεκριμένο πιστοποιητικό και πιστοποιημένη διάταξη δημιουργίας υπογραφής και έχουν έννομο αποτέλεσμα ισοδύναμο με εκείνο της ιδιόχειρης υπογραφής.

\section{Υποδομή δημόσιου κλειδιού και ο κύκλος ζωής των πιστοποιητικών X.509}\label{sec:7.4}
Η κρυπτογραφία δημόσιου κλειδιού θέτει ένα νέο ερώτημα: πώς γνωρίζουμε ότι ένα δημόσιο κλειδί ανήκει πράγματι σε αυτόν που νομίζουμε; Ένας ενδιάμεσος επιτιθέμενος θα μπορούσε να το αντικαταστήσει με το δικό του. Μια \textbf{υποδομή δημόσιου κλειδιού} (ΥΔΚ ή PKI) απαντά σε αυτό το ερώτημα: μια αξιόπιστη \textbf{αρχή πιστοποίησης} (Certificate Authority, CA) υπογράφει ένα \textbf{πιστοποιητικό} που συνδέει ένα δημόσιο κλειδί με μια ταυτότητα, όπως ένα όνομα τομέα. Μια \textbf{αρχή καταχώρισης} (Registration Authority, RA) επαληθεύει την ταυτότητα του αιτούντος πριν η CA εκδώσει το πιστοποιητικό. Τα πιστοποιητικά ακολουθούν το πρότυπο \textbf{X.509} και περιέχουν, μεταξύ άλλων, το υποκείμενο, τον εκδότη, την περίοδο ισχύος, το δημόσιο κλειδί και επεκτάσεις όπως το \emph{Εναλλακτικό Όνομα Υποκειμένου} (Subject Alternative Name, SAN), που απαριθμεί τα ονόματα τομέα που καλύπτει το πιστοποιητικό.

Η εμπιστοσύνη οργανώνεται ως \textbf{αλυσίδα}: μια \emph{ριζική CA} εκτός σύνδεσης, της οποίας το πιστοποιητικό είναι προεγκατεστημένο στα λειτουργικά συστήματα και στους φυλλομετρητές, υπογράφει \emph{ενδιάμεσες CA}, οι οποίες με τη σειρά τους υπογράφουν τα πιστοποιητικά τελικών οντοτήτων. Η επαλήθευση ακολουθεί την αλυσίδα προς τα πάνω έως μια αξιόπιστη ρίζα. Τα πιστοποιητικά έχουν \textbf{κύκλο ζωής} —αίτηση (μέσω αιτήματος υπογραφής πιστοποιητικού, CSR), έκδοση, εγκατάσταση, ανανέωση και, σε περίπτωση διαρροής του κλειδιού, \textbf{ανάκληση}, η οποία γνωστοποιείται μέσω λιστών ανάκλησης (CRL) ή του πρωτοκόλλου άμεσης κατάστασης (OCSP). Επειδή ο έλεγχος ανάκλησης είναι στην πράξη αναξιόπιστος, ο κλάδος κινείται προς πιστοποιητικά μικρής διάρκειας που ανανεώνονται αυτόματα με το πρωτόκολλο ACME.

\section{Το TLS 1.3 και οι εφαρμογές της ΥΔΚ}\label{sec:7.5}
Το \textbf{Transport Layer Security (TLS)} προστατεύει την περιήγηση στον ιστό (HTTPS), τη μεταφορά ηλεκτρονικού ταχυδρομείου, τις διεπαφές προγραμματισμού (API) και πολλά άλλα πρωτόκολλα. Η έκδοση 1.3 (2018) απλοποίησε το πρωτόκολλο και αφαίρεσε παρωχημένες και ευάλωτες επιλογές. Η χειραψία απαιτεί μόνο έναν κύκλο επικοινωνίας. Ο πελάτης στέλνει ένα μήνυμα \emph{ClientHello} με τις σουίτες κρυπτογράφησης που υποστηρίζει και ένα εφήμερο μερίδιο κλειδιού Diffie–Hellman· ο διακομιστής απαντά με το δικό του μερίδιο, το πιστοποιητικό του και μια υπογραφή που αποδεικνύει ότι κατέχει το ιδιωτικό κλειδί. Στη συνέχεια οι δύο πλευρές παράγουν τα ίδια κλειδιά συνεδρίας. Επειδή τα κλειδιά Diffie–Hellman είναι εφήμερα, το TLS 1.3 παρέχει πάντοτε \textbf{εμπρόσθια μυστικότητα} (forward secrecy): η μεταγενέστερη υποκλοπή του ιδιωτικού κλειδιού του διακομιστή δεν επιτρέπει την αποκρυπτογράφηση κίνησης που είχε καταγραφεί στο παρελθόν.

Οι ίδιες αρχές της ΥΔΚ υποστηρίζουν και άλλες εφαρμογές. Το \textbf{S/MIME} υπογράφει και κρυπτογραφεί μηνύματα ηλεκτρονικού ταχυδρομείου· η \textbf{υπογραφή κώδικα} επιτρέπει στα λειτουργικά συστήματα να επαληθεύουν ότι ένα λογισμικό προέρχεται από γνωστό εκδότη και δεν έχει τροποποιηθεί —γι' αυτό τα κλεμμένα πιστοποιητικά που χρησιμοποίησε το Stuxnet ήταν τόσο επιζήμια· και το \textbf{αμοιβαίο TLS (mTLS)} απαιτεί πιστοποιητικό τόσο από τον διακομιστή όσο και από τον πελάτη, πρότυπο που χρησιμοποιείται ευρέως για τον έλεγχο ταυτότητας μεταξύ υπηρεσιών στις σύγχρονες αρχιτεκτονικές.

\section{Διαχείριση κλειδιών: εκεί όπου η κρυπτογραφία πραγματικά αποτυγχάνει}\label{sec:7.6}
Οι σύγχρονοι αλγόριθμοι σπάνια παραβιάζονται μαθηματικά· τα συστήματα αποτυγχάνουν επειδή τα κλειδιά παράγονται πρόχειρα, αποθηκεύονται απρόσεκτα ή δεν ανανεώνονται ποτέ. Τυπικές αστοχίες είναι τα ιδιωτικά κλειδιά που δημοσιεύονται κατά λάθος σε δημόσια αποθετήρια κώδικα, τα ενσωματωμένα διαπιστευτήρια στις εφαρμογές, οι προβλέψιμες γεννήτριες τυχαίων αριθμών και τα πιστοποιητικά που λήγουν χωρίς να το αντιληφθεί κανείς, προκαλώντας διακοπές λειτουργίας. Η \textbf{διαχείριση κλειδιών} καλύπτει ολόκληρο τον κύκλο ζωής: ασφαλή παραγωγή από κρυπτογραφικά ασφαλή πηγή τυχαιότητας, προστατευμένη αποθήκευση, ελεγχόμενη διανομή, τακτική εναλλαγή, ανάκληση και, τέλος, καταστροφή.

Τα κλειδιά υψηλής αξίας πρέπει να φυλάσσονται σε ειδικό υλικό. Μια \textbf{μονάδα ασφάλειας υλικού} (Hardware Security Module, HSM) ή μια υπηρεσία διαχείρισης κλειδιών στο νέφος εκτελεί τις κρυπτογραφικές λειτουργίες εσωτερικά, ώστε το ιδιωτικό κλειδί να μην εγκαταλείπει ποτέ τη συσκευή. Η πρόσβαση στα κλειδιά πρέπει να ακολουθεί τις αρχές του ελάχιστου προνομίου και του διαχωρισμού καθηκόντων, και κάθε χρήση πρέπει να καταγράφεται. Με ματιά στο μέλλον, οι οργανισμοί οφείλουν επίσης να τηρούν απογραφή του πού χρησιμοποιείται κάθε αλγόριθμος (\emph{κρυπτογραφική ευελιξία}), ώστε οι ευάλωτοι αλγόριθμοι να μπορούν να αντικατασταθούν —για παράδειγμα, καθώς υιοθετούνται οι μετακβαντικοί αλγόριθμοι που τυποποίησε το NIST.

\section*{Βασικοί όροι}
\begin{description}[style=nextline,leftmargin=1.2em]
  \item[Υβριδική κρυπτογράφηση] Χρήση ασύμμετρης κρυπτογραφίας για την προστασία ενός συμμετρικού κλειδιού συνεδρίας που κρυπτογραφεί τα δεδομένα.
  \item[Αντίσταση συγκρούσεων] Η αδυναμία εύρεσης δύο διαφορετικών εισόδων με την ίδια σύνοψη.
  \item[Ψηφιακή υπογραφή] Τιμή που υπολογίζεται με ιδιωτικό κλειδί και αποδεικνύει την ακεραιότητα και την προέλευση ενός μηνύματος.
  \item[Αρχή πιστοποίησης] Αξιόπιστη οντότητα που υπογράφει πιστοποιητικά τα οποία συνδέουν δημόσια κλειδιά με ταυτότητες.
  \item[Εμπρόσθια μυστικότητα] Η ιδιότητα κατά την οποία η διαρροή μακροπρόθεσμων κλειδιών δεν εκθέτει την κίνηση προηγούμενων συνεδριών.
  \item[HSM] Μονάδα ασφάλειας υλικού: ανθεκτική σε παραβίαση συσκευή που φυλάσσει κλειδιά και εκτελεί κρυπτογραφικές λειτουργίες.
\end{description}

\section*{Σύνοψη κεφαλαίου}
\begin{itemize}
  \item Οι συμμετρικοί αλγόριθμοι είναι γρήγοροι αλλά απαιτούν κοινό κλειδί· οι ασύμμετροι επιλύουν τη διανομή κλειδιών· τα υβριδικά σχήματα συνδυάζουν και τα δύο.
  \item Οι συναρτήσεις κατακερματισμού παρέχουν έλεγχο ακεραιότητας και όχι εμπιστευτικότητα· οι κωδικοί απαιτούν αργές συναρτήσεις κατακερματισμού με αλάτι.
  \item Οι ψηφιακές υπογραφές παρέχουν ακεραιότητα και αυθεντικότητα· η μη αποποίηση εξαρτάται επιπλέον από τη φύλαξη των κλειδιών και το νομικό πλαίσιο.
  \item Η ΥΔΚ συνδέει κλειδιά με ταυτότητες μέσω αλυσίδων πιστοποιητικών· η διαχείριση του κύκλου ζωής και η ανάκληση είναι απαραίτητες.
  \item Το TLS 1.3 προσφέρει γρήγορη χειραψία με εμπρόσθια μυστικότητα, όμως οι περισσότερες πραγματικές αποτυχίες οφείλονται σε κακή διαχείριση κλειδιών.
\end{itemize}

\section*{Ερωτήσεις ανασκόπησης}
\begin{enumerate}
  \item Γιατί τα πρακτικά συστήματα δεν κρυπτογραφούν μεγάλα αρχεία απευθείας με RSA;
  \item Ένας προγραμματιστής αποθηκεύει τους κωδικούς ως συνόψεις SHA-256 χωρίς αλάτι. Εξηγήστε τον κίνδυνο και προτείνετε καλύτερο σχεδιασμό.
  \item Περιγράψτε βήμα προς βήμα πώς ένας φυλλομετρητής επαληθεύει την αλυσίδα πιστοποιητικών ενός ιστοτόπου.
  \item Εξηγήστε την εμπρόσθια μυστικότητα και γιατί το TLS 1.3 κατάργησε τη στατική ανταλλαγή κλειδιών RSA.
\end{enumerate}
